%! Logarithmic Function Context and Data Modeling — Unit 2, Lesson 2.14 — Group Activity
\documentclass[10pt]{article}
\usepackage{precalculus-article}
\usepackage{precalculus-boxes}
\newcommand{\TallMath}[1]{$\displaystyle #1\rule[-1.4em]{0pt}{3.2em}$}

\begin{document}

\pageheader{Unit 2, Lesson 2.14}{Group Activity}
\namepartnerperiod

\vspace{0.1in}

\begin{scenariobox}[Species Diversity and Sound: Logarithmic Models in Action]{sky}
An ecologist studying habitat biodiversity and a physicist studying sound both encounter
data that are best modeled with logarithmic functions. Your group will construct, interpret,
and apply logarithmic models in each setting.
\end{scenariobox}

\vspace{0.1in}

\begin{tcolorbox}[colback=white, colframe=black!40,
    title=\textbf{Tier R --- Remediate},
    fonttitle=\bfseries, arc=2mm, left=3mm, right=3mm, top=2mm, bottom=2mm]

\textbf{Species Diversity --- reading a table and using a simple model.}

\vspace{4pt}
The table below shows how species count $S$ depends on habitat area $A$ (hectares).

\vspace{4pt}
\begin{center}
\renewcommand{\arraystretch}{1.8}
\begin{tabular}{|c|c|c|c|c|}
\hline
\rowcolor{sky}
$A$ & $1$ & $10$ & $100$ & $1000$ \\
\hline
$S$ & $0$ & $1$  & $2$   & $3$ \\
\hline
\end{tabular}
\end{center}

\begin{enumerate}[leftmargin=1.8em, itemsep=10pt]

\item As $S$ increases by 1, how does $A$ change? (Describe the pattern.)

\vspace{4pt}
$A$ \blank{5cm}. The proportion factor is: \blank{1.5cm}

\item Using the proportion and the fact that $S = 0$ when $A = 1$ (the real zero), write
      the logarithmic model.

\vspace{4pt}
$S(A) = $ \blank{4cm}

\item Use the model to predict $S(10000)$.

\vspace{4pt}
$S(10000) = $ \blank{5cm}

\item Is $S(0.5)$ defined? Explain why or why not in terms of the domain of a
      logarithmic function.

\writelines{2}

\end{enumerate}
\end{tcolorbox}

\vspace{0.1in}

\begin{tcolorbox}[colback=white, colframe=black!40,
    title=\textbf{Tier A --- Approaching Proficiency},
    fonttitle=\bfseries, arc=2mm, left=3mm, right=3mm, top=2mm, bottom=2mm]

\textbf{Species Diversity --- building a model from two data points.}

\vspace{4pt}
A different study gives two data points for species count $S$ as a function of area $A$:
when $A = 1$, $S = 2$; and when $A = 9$, $S = 4$. Assume the model has the form
$S(A) = a\log_3(A) + k$.

\begin{enumerate}[leftmargin=1.8em, itemsep=10pt]

\item Write the system of two equations using the two data points.

\vspace{4pt}
Equation 1 (use $A=1$, $S=2$): \blank{6cm}

\vspace{4pt}
Equation 2 (use $A=9$, $S=4$): \blank{6cm}

\item Since $\log_3(1) = $ \blank{1cm}, use Equation 1 to find $k$.

\vspace{4pt}
$k = $ \blank{2cm}

\item Substitute $k$ into Equation 2 to find $a$. Show your work.

\vspace{4pt}
\writelines{2}
$a = $ \blank{2cm}

\item Write the complete model.

\vspace{4pt}
$S(A) = $ \blank{5cm}

\item Predict the species count when $A = 27$.

\vspace{4pt}
$S(27) = $ \blank{6cm}

\end{enumerate}
\end{tcolorbox}

\vspace{0.1in}

\begin{tcolorbox}[colback=white, colframe=black!40,
    title=\textbf{Tier E --- Extension},
    fonttitle=\bfseries, arc=2mm, left=3mm, right=3mm, top=2mm, bottom=2mm]

\textbf{Sound intensity --- transformations and real-world predictions.}

\vspace{4pt}
The decibel formula is $S(I) = 10\log_{10}\!\left(\dfrac{I}{I_0}\right)$,
where $I$ is sound intensity in W/m$^2$ and $I_0 = 10^{-12}$ W/m$^2$ is the
threshold of hearing.

\begin{enumerate}[leftmargin=1.8em, itemsep=10pt]

\item Use the property $\log_{10}(A/B) = \log_{10}(A) - \log_{10}(B)$ to expand:

\vspace{4pt}
$S(I) = 10\log_{10}(I) - 10\log_{10}(I_0) = 10\log_{10}(I) + $ \blank{3cm}

\vspace{2pt}
(Hint: $\log_{10}(10^{-12}) = $ \blank{2cm}, so $-10\log_{10}(I_0) = $ \blank{2cm}.)

\item Identify $S(I) = 10\log_{10}(I) + 120$ as a transformation of $f(I) = 10\log_{10}(I)$.
      What is the vertical shift?

\vspace{4pt}
Vertical shift: \blank{3cm}

\item Normal conversation is 60 dB. Find the intensity $I$ for $S = 60$.

\vspace{4pt}
$60 = 10\log_{10}(I) + 120$

\vspace{4pt}
$\log_{10}(I) = $ \blank{2cm}

\vspace{4pt}
$I = $ \blank{4cm} W/m$^2$

\item A rock concert is 120 dB. Find the intensity for $S = 120$, then compute how many
      times more intense a rock concert is than normal conversation.

\vspace{4pt}
Intensity at 120 dB: $I = $ \blank{4cm} W/m$^2$

\vspace{4pt}
Ratio: \blank{6cm} times more intense

\end{enumerate}
\end{tcolorbox}

\end{document}
